\clearpage
\begin{explanationbox}
\textbf{\textcolor{CExplanation}{一、为什么要有这条结论？}}
\begin{center}\begin{tikzpicture}[font=\small,>=stealth,line join=round]
\draw[gray!40,rounded corners=3pt] (0,0) rectangle (14.2,8.3);\node at (7.1,7.700000000000001) {镜子在1与3：半格出发，两次反射后到四点五};
\begin{scope}[shift={(1.6901639344262296,3.814705882352941)},x=0.9836065573770493cm,y=1.352941176470588cm]\draw[->,gray] (-0.6,0)--(5.5,0) node[right]{$x$};\draw[->,gray] (0,-1.6)--(0,1.8) node[above]{$y$};\draw[blue,thick,solid,domain=-0.2:5,samples=180] plot(\x,{cos(deg(pi*(\x-1)/2))});\fill[orange!85!black] (0.5,0.7071067811865476) circle(2.2pt);\fill[orange!85!black] (1.5,0.7071067811865476) circle(2.2pt);\fill[orange!85!black] (4.5,0.7071067811865476) circle(2.2pt);\foreach \a in {1,3}{\draw[violet,dashed] (\a,-1.3)--(\a,1.3);\node[above] at (\a,1.3){$x=\a$};}\draw[orange!85!black,dashed] (.5,.70710678)--(4.5,.70710678);\end{scope}\node[anchor=west,align=left,text width=5.0cm] at (8.4,4.0) {$f=\cos\frac{\pi(x-1)}2$\\[9pt]轴 $x=1,3$\\[9pt]$0.5\to1.5\to4.5$\\[9pt]三点同高 $\frac{\sqrt2}2$\\[9pt]横移4，两倍轴间距};
\node[text width=12.6cm,align=center] at (7.1,.65) {图中反射点来自曲线；任意起点同样满足两次反射后水平移动4。};
\end{tikzpicture}\end{center}



遇到 $f(1+x)=f(1-x)$ 时，我们知道图像关于 $x=1$ 对称；再遇到关于 $x=3$ 对称时，仅凭两条信息，竟然就能知道：\textbf{整条图像每向右移动 $4$ 个单位就重复一次}。
这让我们在\textbf{不知道 $f(x)$ 具体是什么函数}时，也能求周期、算遥远位置的函数值。

\medskip

\end{explanationbox}
\clearpage
\begin{explanationbox}
\textbf{\textcolor{CExplanation}{二、把“对称”想成照镜子：先从初中能理解的数开始}}
\begin{center}\begin{tikzpicture}[font=\small,>=stealth,line join=round]
\draw[gray!40,rounded corners=3pt] (0,0) rectangle (14.2,8.3);\node at (7.1,7.700000000000001) {例2：同高中心翻两次，5变成1，再回到5};
\begin{scope}[shift={(2.4846153846153847,1.9785714285714284)},x=0.7692307692307693cm,y=0.6571428571428571cm]\draw[->,gray] (-1.8,0)--(6,0) node[right]{$x$};\draw[->,gray] (0,-0.5)--(0,6.5) node[above]{$y$};\draw[teal,thick,solid,domain=-1.4:5.6,samples=180] plot(\x,{3+2*sin(deg(pi*\x/2))});\fill[orange!85!black] (1,5) circle(2.2pt);\fill[orange!85!black] (-1,1) circle(2.2pt);\fill[orange!85!black] (5,5) circle(2.2pt);\fill[orange!85!black] (0,3) circle(2.2pt);\fill[orange!85!black] (2,3) circle(2.2pt);\draw[gray,dashed] (-1.6,3)--(5.8,3);\end{scope}\node[anchor=west,align=left,text width=5.0cm] at (8.4,4.0) {$f=3+2\sin\frac{\pi x}2$\\[9pt]中心 $(0,3),(2,3)$\\[9pt]$1\to-1\to5$\\[9pt]高度 $5\to1\to5$\\[9pt]正周期可取4};
\node[text width=12.6cm,align=center] at (7.1,.65) {两个中心共用高度3，纵坐标经过两次关于3的翻转后还原。};
\end{tikzpicture}\end{center}



\textbf{1. 两条镜子为什么是两倍间距？}

假设镜子在 $x=1$ 和 $x=3$，两镜子相距 $2$。从位置 $x=0.5$ 出发：
\[
\underbrace{0.5}_{\text{原位置}}
\xrightarrow[\text{照 }x=1]{\text{左右翻}}
1.5
\xrightarrow[\text{照 }x=3]{\text{左右翻}}
4.5.
\]

两次翻折，$0.5$ 变为 $4.5$，正好向右移了 $4=2\times2$。
对\textbf{任何}起点 $x$ 都一样：
\[
 x\longrightarrow 2a-x\longrightarrow 2b-(2a-x)=x+2(b-a).
\]
两个箭头仅改变位置，不改变函数值，所以 $f(x+2(b-a))=f(x)$。

\textbf{2. 绕一个点转半圈，为什么还要管高度？}

若对称中心是 $(b,c)$，横坐标从 $x$ 变成 $2b-x$，同时纵坐标 $y$ 会变成 $2c-y$。
比如中心在高度 $c=3$：$5$ 在中线上方 $2$ 格，转半圈后变为 $1$（中线下方 $2$ 格）：
\[
 \underbrace{5}_{\text{上方2格}}\longrightarrow \underbrace{2\times3-5=1}_{\text{下方2格}}
 \longrightarrow \underbrace{2\times3-1=5}_{\text{回到上方2格}}.
\]

因此，如果\textbf{两个点都位于 $y=c$}，连续绕这两个点转半圈，横向共移动 $2d$，纵向也回到原来的高度：点点组合得到 $T=2d$。

\end{explanationbox}\clearpage\begin{explanationbox}
\begin{center}\begin{tikzpicture}[font=\small,>=stealth,line join=round]
\draw[gray!40,rounded corners=3pt] (0,0) rectangle (14.2,8.3);\node at (7.1,7.700000000000001) {例3：轴在0，中心 $(1,3)$；走2高度尚未还原};
\begin{scope}[shift={(2.174626865671642,1.9785714285714284)},x=0.8955223880597015cm,y=0.6571428571428571cm]\draw[->,gray] (-1.2,0)--(5.5,0) node[right]{$x$};\draw[->,gray] (0,-0.5)--(0,6.5) node[above]{$y$};\draw[blue,thick,solid,domain=-0.8:5,samples=180] plot(\x,{3+2*cos(deg(pi*\x/2))});\fill[orange!85!black] (0,5) circle(2.2pt);\fill[orange!85!black] (2,1) circle(2.2pt);\fill[orange!85!black] (4,5) circle(2.2pt);\fill[orange!85!black] (1,3) circle(2.2pt);\draw[violet,dashed] (0,-.3)--(0,6);\draw[gray,dashed] (-1,3)--(5.2,3);\end{scope}\node[anchor=west,align=left,text width=5.0cm] at (8.4,4.0) {$f=3+2\cos\frac{\pi x}2$\\[9pt]轴 $x=0$，中心 $(1,3)$\\[9pt]$f(0)=5,\ f(2)=1$\\[9pt]$f(4)=5$\\[9pt]$d=1,\ T=4$};
\node[text width=12.6cm,align=center] at (7.1,.65) {相对高度3一高一低；移动2只翻转，移动4才保证原样重复。};
\end{tikzpicture}\end{center}

\textbf{3. 一面镜子加一个转弯点，为什么要四倍？}

镜子只把左、右交换，\textbf{不把高、低交换}；中心对称同时交换左右与高低。
假设 $a=0,b=1,c=3$，则两处相距 $d=1$：
\[
\begin{aligned}
  f(x+2)&=6-f(x),\\
  f(x+4)&=6-f(x+2)=f(x).
\end{aligned}
\]
例如 $f(0)=5$，就必须有 $f(2)=1$、$f(4)=5$。
\textbf{走 $2$ 格只是“上下相反”，走 $4$ 格才“原样回来”。}

\medskip

\end{explanationbox}
\clearpage
\begin{explanationbox}
\textbf{\textcolor{CExplanation}{三、三张图：亲眼看见两倍和四倍}}
\begin{center}\begin{tikzpicture}[font=\small,>=stealth,line join=round]
\draw[gray!40,rounded corners=3pt] (0,0) rectangle (14.2,8.5);\node at (7.1,7.9) {同样连续对称：两次后是否回到原高度？};
\node at (2.4,6.9) {轴轴：$f=\cos\frac{\pi x}2$};\begin{scope}[shift={(1.2666666666666666,4.348571428571429)},x=0.5666666666666667cm,y=1.0285714285714287cm]\draw[->,gray] (-1,0)--(5,0) node[right]{$x$};\draw[->,gray] (0,-1.7)--(0,1.8) node[above]{$y$};\draw[blue,thick,solid,domain=-0.7:4.8,samples=160] plot(\x,{cos(deg(pi*\x/2))});\foreach \a in {0,2}{\draw[violet,dashed] (\a,-1.4)--(\a,1.4);}\end{scope}\node[align=center,text width=4cm] at (2.4,1.7) {轴 $x=0,2$\\$d=2,\ T=4$};\node at (7.050000000000001,6.9) {点点：$g=\sin\frac{\pi x}2$};\begin{scope}[shift={(5.916666666666667,4.348571428571429)},x=0.5666666666666667cm,y=1.0285714285714287cm]\draw[->,gray] (-1,0)--(5,0) node[right]{$x$};\draw[->,gray] (0,-1.7)--(0,1.8) node[above]{$y$};\draw[teal,thick,solid,domain=-0.7:4.8,samples=160] plot(\x,{sin(deg(pi*\x/2))});\fill[orange!85!black] (0,0) circle(2pt);\fill[orange!85!black] (2,0) circle(2pt);\end{scope}\node[align=center,text width=4cm] at (7.050000000000001,1.7) {中心 $(0,0),(2,0)$\\$d=2,\ T=4$};\node at (11.7,6.9) {轴点：同一条余弦线};\begin{scope}[shift={(10.566666666666666,4.348571428571429)},x=0.5666666666666667cm,y=1.0285714285714287cm]\draw[->,gray] (-1,0)--(5,0) node[right]{$x$};\draw[->,gray] (0,-1.7)--(0,1.8) node[above]{$y$};\draw[blue,thick,solid,domain=-0.7:4.8,samples=160] plot(\x,{cos(deg(pi*\x/2))});\fill[orange!85!black] (1,0) circle(2pt);\draw[violet,dashed] (0,-1.4)--(0,1.4);\end{scope}\node[align=center,text width=4cm] at (11.7,1.7) {轴0，中心 $(1,0)$\\$d=1,\ T=4$};
\node[text width=12.6cm,align=center] at (7.1,.65) {三图同尺度：轴轴、同高点点得到2d；轴点半程翻转高度，需4d。};
\end{tikzpicture}\end{center}



三图同尺度直接比较：同一条余弦线选轴0与轴2，或选轴0与中心1，都会得到周期4；区别在选取的对称位置和距离。


\end{explanationbox}
\clearpage
\begin{explanationbox}
\textbf{\textcolor{CExplanation}{四、和已经学过的知识连起来}}
\begin{center}\begin{tikzpicture}[font=\small,>=stealth,line join=round]
\draw[gray!40,rounded corners=3pt] (0,0) rectangle (14.2,8.3);\node at (7.1,7.700000000000001) {证明三：整条曲线下移3，中心高度变成0};
\begin{scope}[shift={(2.207692307692308,3.102631578947368)},x=0.9230769230769231cm,y=0.4842105263157894cm]\draw[->,gray] (-1.2,0)--(5.3,0) node[right]{$x$};\draw[->,gray] (0,-3)--(0,6.5) node[above]{$y$};\draw[blue,thick,solid,domain=-0.8:4.8,samples=180] plot(\x,{3+2*cos(deg(pi*\x/2))});\draw[orange!85!black,thick,dashed,domain=-0.8:4.8,samples=180] plot(\x,{2*cos(deg(pi*\x/2))});\fill[orange!85!black] (1,3) circle(2.2pt);\fill[orange!85!black] (1,0) circle(2.2pt);\draw[gray,dashed] (-1,3)--(5,3);\draw[<->,violet] (1.2,0)--(1.2,3);\end{scope}\node[anchor=west,align=left,text width=5.0cm] at (8.4,4.0) {蓝：$f=3+2\cos\frac{\pi x}2$\\[9pt]橙虚线：$g=f-3$\\[9pt]中心 $(1,3)\to(1,0)$\\[9pt]周期都为4\\[9pt]纵向平移不改重复距离};
\node[text width=12.6cm,align=center] at (7.1,.65) {把中线当作新的零线，关于中心的高度翻转就变成正负号交换。};
\end{tikzpicture}\end{center}


\begin{itemize}[leftmargin=*,itemsep=0.35em]
\item \textbf{与偶函数相连：}关于 $x=a$ 对称，就是把原点平移到 $a$ 后，左右两边函数值相同。特别地，$a=0$ 时是 $f(x)=f(-x)$。
\item \textbf{与奇函数相连：}关于 $(b,c)$ 对称，就是先把纵向基准线移到 $y=0$，再把横向原点移到 $b$，左右函数值相反。特别地，$(b,c)=(0,0)$ 时是 $f(-x)=-f(x)$。
\item \textbf{与周期定义相连：}只有当\textbf{所有} $x$ 都满足 $f(x+T)=f(x)$，才能称 $T>0$ 为周期；只核对图上的几个点不能当作证明。
\end{itemize}
\end{explanationbox}
